
\chapter{Überleben I -- Die Begrüßungs-Choreografie}

In Ägypten ist eine Begrüßung kein einzelnes Wort -- es ist ein kleines
Ritual mit festen Schritten, wie ein Tanz. Wenn du den ersten Schritt
machst, \emph{musst} die andere Person mit dem passenden zweiten Schritt
antworten. Das ist keine Höflichkeit, das ist Choreografie.

\begin{itemize}[leftmargin=*]
\item \arword{السلام عليكم}{as-salámu 3aláykum}{Der Friede sei mit dir (universeller Gruß)}{\%D8\%A7\%D9\%84\%D8\%B3\%D9\%84\%D8\%A7\%D9\%85\%20\%D8\%B9\%D9\%84\%D9\%8A\%D9\%83\%D9\%85}
\item \arword{وعليكم السلام}{wa 3aláykum as-saláam}{Antwort: und mit dir sei der Friede}{\%D9\%88\%D8\%B9\%D9\%84\%D9\%8A\%D9\%83\%D9\%85\%20\%D8\%A7\%D9\%84\%D8\%B3\%D9\%84\%D8\%A7\%D9\%85}
\item \arword{صباح الخير}{SabáaH el-kheir}{Guten Morgen}{\%D8\%B5\%D8\%A8\%D8\%A7\%D8\%AD\%20\%D8\%A7\%D9\%84\%D8\%AE\%D9\%8A\%D8\%B1}
\item \arword{صباح النور}{SabáaH en-nuur}{Antwort auf Guten Morgen}{\%D8\%B5\%D8\%A8\%D8\%A7\%D8\%AD\%20\%D8\%A7\%D9\%84\%D9\%86\%D9\%88\%D8\%B1}
\item \arword{مساء الخير}{masáa' el-kheir}{Guten Abend}{\%D9\%85\%D8\%B3\%D8\%A7\%D8\%A1\%20\%D8\%A7\%D9\%84\%D8\%AE\%D9\%8A\%D8\%B1}
\item \arword{مساء النور}{masáa' en-nuur}{Antwort auf Guten Abend}{\%D9\%85\%D8\%B3\%D8\%A7\%D8\%A1\%20\%D8\%A7\%D9\%84\%D9\%86\%D9\%88\%D8\%B1}
\item \arword{إزيك}{izzáyyak}{Wie geht's dir? (zu einem Mann, Masri)}{\%D8\%A5\%D8\%B2\%D9\%8A\%D9\%83}
\item \arword{إزيك}{izzáyyik}{Wie geht's dir? (zu einer Frau, gleiche Schrift!)}{\%D8\%A5\%D8\%B2\%D9\%8A\%D9\%83}
\item \arword{كويس، الحمد لله}{kwáyyes, el-Hamdu lilláh}{Gut, Gott sei Dank (Standardantwort -- immer!)}{\%D9\%83\%D9\%88\%D9\%8A\%D8\%B3\%D8\%8C\%20\%D8\%A7\%D9\%84\%D8\%AD\%D9\%85\%D8\%AF\%20\%D9\%84\%D9\%84\%D9\%87}
\item \arword{شكراً}{shukran}{Danke}{\%D8\%B4\%D9\%83\%D8\%B1\%D8\%A7\%D9\%8B}
\item \arword{العفو}{el-3afw}{Bitte/Gerne (Antwort auf Danke)}{\%D8\%A7\%D9\%84\%D8\%B9\%D9\%81\%D9\%88}
\item \arword{من فضلك}{min faDlak}{Bitte (zu einem Mann)}{\%D9\%85\%D9\%86\%20\%D9\%81\%D8\%B6\%D9\%84\%D9\%83}
\item \arword{من فضلك}{min faDlik}{Bitte (zu einer Frau)}{\%D9\%85\%D9\%86\%20\%D9\%81\%D8\%B6\%D9\%84\%D9\%83}
\item \arword{مع السلامة}{ma3a s-saláama}{Tschüss (,,geh in Frieden``)}{\%D9\%85\%D8\%B9\%20\%D8\%A7\%D9\%84\%D8\%B3\%D9\%84\%D8\%A7\%D9\%85\%D8\%A9}
\item \arword{تصبح على خير}{tuSbaH 3ala kheir}{Gute Nacht (zu einem Mann)}{\%D8\%AA\%D8\%B5\%D8\%A8\%D8\%AD\%20\%D8\%B9\%D9\%84\%D9\%89\%20\%D8\%AE\%D9\%8A\%D8\%B1}
\end{itemize}

\begin{ammahmoudbox}
,,Egal wie dein Tag war -- auf 'izzayyak?' antwortest du IMMER mit
'kwayyes, el-Hamdu lillah'. Auch wenn dein Taxi gerade eine Panne hatte.
Das ist keine Lüge, das ist Etikette.``
\end{ammahmoudbox}

\begin{lachbox}
Die Antwort auf ,,Wie geht's?`` in Ägypten ist wie das deutsche
,,gut, und dir?`` -- nur, dass ,,gut`` hier gesetzlich vorgeschrieben zu
sein scheint. Niemand antwortet ehrlich. Das ist Teil des Charmes.
\end{lachbox}

\begin{checkpointbox}
Übe laut (am besten mit \textcolor{yallateal}{\faVolumeUp}): ,,as-sal\'amu
3al\'aykum`` sagen, kurze Pause, dann selbst mit ,,wa 3al\'aykum
as-sal\'aam`` antworten. Wenn das flüssig klingt, bist du bereit für echte
Menschen.
\end{checkpointbox}
